\section{Introduction and statement}
\label{sec:introduction}

The minimal model program seeks to simplify a projective variety by operations
on which its canonical divisor, or an adjoint divisor $K_X+B$, is negative.
A divisorial contraction removes a divisor. A flip replaces a small
contraction by a birational model on which the adjoint becomes relatively
ample. The expected endpoint is a nef adjoint or a Mori fibre space.
Existence of the required operations makes the program possible;
termination ensures that a sequence of choices actually reaches an endpoint.
It is stronger than the existence of one minimal model or the termination
of a particular program with scaling.

The three-dimensional log minimal model program established the basic
pattern; see \cite{KollarEtAl1992,KM98}. In higher dimensions,
Birkar--Cascini--Hacon--McKernan proved klt flip existence and
minimal-model existence for klt pairs with big boundary and
pseudo-effective adjoint \cite{BCHM}.
Fujino's cone and contraction theorem and Birkar's lc flip theorem
supply the corresponding operations for log canonical pairs
\cite{FujinoFoundations,BirkarLC}. In dimension four, Shokurov's
ordered-termination method constructed minimal models or Mori fibre
spaces; Birkar gave a short proof of the klt case
\cite{ShokurovOrdered,BirkarShokurov}. These constructions use
selected programs, and do not establish termination for every
successive choice of negative ray.

The difficulty method studies arbitrary sequences by combining
discrepancy improvement with the geometry of the exceptional loci.
Shokurov introduced discrepancy-counting difficulty
\cite[Definition~2.15 and Corollaries~2.16--2.17]{Shokurov1985}.
Kawamata proved termination for terminal fourfolds
\cite{KMM87}, and Fujino extended the method to
canonical pairs with rational boundary; the weighted difficulty
and precise hypotheses are corrected in
\cite{FujinoCanonical,FujinoCanonicalAddendum}.
Algebraic cycle classes account for the remaining flips of
type $(2,1)$ in these arguments. Alexeev--Hacon--Kawamata developed
weighted difficulties that subtract the families of valuations
arising along a boundary component and include cycle-rank terms
on its normalization \cite{AHK}. Their Theorem~3.2 proves
termination of klt fourfold flips when $-(K_X+B)$ is numerically equivalent to an
effective divisor. This numerical effectivity hypothesis does
not follow from non-pseudo-effectivity of $K_X+B$.

Arbitrary termination for pseudo-effective lc fourfolds is now
available through the generalized-pair theorem of
Chen--Tsakanikas \cite[Theorem~1.1]{ChenTsakanikas} and Moraga's
termination theorem \cite[Theorem~1]{Moraga2025}. Both cover
ordinary pairs with real boundary and impose no initial
$\mathbb Q$-factoriality in their statements about sequences of
flips. The unrestricted theorem of Chen--Tsakanikas is
three-dimensional. More recently, Han--Liu--Zhuang proved
termination of every klt fourfold MMP with big boundary, including
the non-pseudo-effective case
\cite[Corollary~3.2(3)]{HanLiuZhuang2025}.

The result below treats ordinary lc fourfolds with rational
boundary without a pseudo-effectivity or big-boundary hypothesis.
It also includes mixed birational operations on models which
need not be $\mathbb Q$-factorial. The proof follows the difficulty
approach: its additional ingredient controls the normalization
branches that enter the correction terms when coefficients on
the auxiliary terminal models continue to decrease.

We prove the projective four-dimensional case of the Termination of Flips
Conjecture for ordinary log canonical pairs with rational boundary;
see \cite[Introduction]{ChenTsakanikas} for the conjecture and its
generalized-pair formulation. Our program always starts from the given
model. We first prove the theorem over $\mathbb C$, where the branch
estimate uses topology, and then transfer the full assertion to every
algebraically closed characteristic-zero field.

Throughout, a pair $(X,B)$ consists of a normal integral variety and an effective $\mathbb Q$-divisor such that $K_X+B$ is $\mathbb Q$-Cartier. Canonical divisors on birational models are chosen using the same rational top differential. For a normalized divisorial valuation $v=\operatorname{ord}_E$, with $E$ on a smooth model $p:W\to X$, our convention is
\[
 a(v;X,B)=1+\operatorname{coeff}_E\bigl(K_W-p^*(K_X+B)\bigr).
\]
Thus $a$ is the \emph{log discrepancy}. Its center $c_X(v)$ is the closed irreducible center. A valuation is exceptional over $X$ when this center has codimension at least two. Log canonical (lc) and klt mean, respectively, that every log discrepancy is nonnegative and positive. A terminal pair is klt and has $a(v;X,B)>1$ for every exceptional valuation. The boundaries in the pair statements are effective; the crepant pullback divisors used on log resolutions may have negative coefficients.

\begin{definition}[Programs and mixed steps]\label{def:program}
Let $(X,B)$ be a projective lc pair. Set $(X_0,B_0)=(X,B)$ and $D_i=K_{X_i}+B_i$. At a non-nef stage choose any $D_i$-negative extremal ray of $\overline{\operatorname{NE}}(X_i)$ and its projective contraction
\[
 f_i:X_i\longrightarrow Z_i,
 \qquad (f_i)_*\mathcal O_{X_i}=\mathcal O_{Z_i},
 \qquad \rho(X_i/Z_i)=1.
\]
The target is normal, and $-D_i$ is relatively ample. Stop if $f_i$ is of fibre type. If it is birational, a permitted step is
\[
 X_i\xrightarrow{\ f_i\ }Z_i
 \xleftarrow{\ f_i^+\ }X_{i+1},
\]
where $f_i^+$ is small and projective birational, possibly an isomorphism, $(X_{i+1},B_{i+1})$ is lc, $B_{i+1}$ is the strict transform of $(f_i)_*B_i$, and $D_{i+1}$ is $\mathbb Q$-Cartier and $f_i^+$-ample. There is no condition on $\rho(X_{i+1}/Z_i)$.

This includes divisorial contractions, flips, and a divisorial contraction followed by a small modification. We call the last possibility a \emph{mixed step}. Stop at a nef adjoint as well. Flops, extra blowups, and inserted isomorphisms are not steps. A program is maximal if it is continued whenever neither stopping condition holds. The broader face-contracting convention in \cite[Definition~2.5]{HanLiuZhuang2025} also allows such birational diagrams; here the negative contraction is required to contract one extremal ray.
\end{definition}

\begin{theorem}\label{thm:main}
Let $k$ be an algebraically closed field of characteristic zero, and let $(X,B)$ be a projective lc pair of dimension four with rational boundary. No $\mathbb Q$-factoriality or pseudo-effectivity is assumed.
\begin{enumerate}[label=\textup{(\roman*)}]
\item Every negative extremal ray at every stage has the contraction and, when birational, a positive model required by Definition~\ref{def:program}.
\item Every program of Definition~\ref{def:program} starting from $(X,B)$ has finitely many birational steps, for every sequence of permitted choices.
\item Every maximal program ends either at a projective lc pair with nef adjoint, or at a Mori fibre space: a projective contraction with connected fibres, relative Picard number one, lower-dimensional normal base, and relatively ample negative adjoint.
\end{enumerate}
An already nef input has a zero-step completed program.
\end{theorem}

The auxiliary terminal and dlt models used in the proof retain the
initial model and every prescribed downstairs choice.  Theorem~\ref{thm:main}
is a termination statement; semiampleness of a nef endpoint requires
abundance.  Section~\ref{sec:endpoints} gives that consequence using
\cite{OpenAILogAbundance2026}.  It also obtains uniform local complements
and Cartier sections at $\epsilon$-lc Mori endpoints over $\mathbb C$
from \cite{OpenAIComplements2026} and \cite{OpenAICartier2026}, respectively.

\subsection*{The difficulty and the branch estimate}

The proof reduces termination to a numerical problem on auxiliary
terminal models.  Ranks of divisor cycle classes first remove all
steps that lose divisors, including mixed steps.  A remaining klt
small tail lifts to crepant terminal junctions, joined by coefficient
decreases and finite small terminal programs.  The same finite set
of boundary labels persists throughout this chain.  Surface minimal
log discrepancy ACC confines every label whose coefficient keeps
varying to a divisor mapped downstairs to a curve or a point.

A weighted discrepancy count on these models has an ordinary family
of contributions at each surface on a single boundary component.
We subtract that family's default weight locally, before summing,
and add divisor-cycle ranks on the boundary normalizations.  This
produces a finite integer $\mathcal D$.  The weights are ceilings
\(w(u)=\lceil N\max\{u,0\}\rceil\), with $N$ clearing only the
original and fixed coefficients; a valuation of log discrepancy $a$
has weight $w(2-a)$.  At a surface $V$ with $r\geq1$ varying
normalization branches, the sum of their individual defaults can
exceed the weight paid by the first blowup.  The elementary inequality
\[
 \lceil x_1+\cdots+x_r\rceil
 \geq\sum_{j=1}^r\lceil x_j\rceil-(r-1)
\]
identifies the possible deficit as $r-1$.  Section~\ref{sec:difficulty-setup}
constructs $\mathcal D$ and proves this local bound, including the
higher default contributions.

The geometric input, Theorem~\ref{thm:branch}, bounds the total excess
of normalization branches along surfaces of the varying boundary.
It pays for that excess using divisor-cycle ranks of its normalized
components and distinct valuations centered on curves, with log
discrepancy below two.  These are already positive terms in
$\mathcal D$, so the estimate makes $\mathcal D$ nonnegative at every
crepant junction.  The branches are algebraic prime divisors over
the generic field of the surface; residue degrees enter their
cycle pushforward.

To prove the estimate, normalization relations first give a bound
through weight-four degree-five cohomology.  Projective decomposition
and a support estimate isolate punctured degree-two local systems
along singular curves.  Finite monodromy and norms represent invariant
classes by line bundles over each curve's original function field.
Formal existence and approximation extend these bundles to a pointed \'etale
neighbourhood with unchanged residue field.  Degrees on a small
model produce distinct divisorial valuations of the original function
field, each with zero-boundary log discrepancy two.  The positive
boundary lowers their discrepancies strictly.  The topological and
valuation arguments occupy Sections~\ref{sec:branch-topology}
and~\ref{sec:local-valuations}.

Section~\ref{sec:difficulty} returns to the terminal chain.  Exact
cancellation between normalization ranks and branch defaults makes
$\mathcal D$ nonincreasing.  Every positive-side bad surface forces
a unit decrease: the relevant endpoint discrepancy has an integral
scaled deficit, so strict discrepancy improvement crosses a ceiling
threshold.  Hence only finitely many such surfaces occur.  A second
cycle-rank comparison and the graph dimension bound finish klt
termination.  Section~\ref{sec:lc} uses special termination and
nonempty dlt bridges to transfer this conclusion to each prescribed
lc program, and proves continuation through every chosen negative
ray.  Section~\ref{sec:field} transfers an entire hypothetical
countable program to $\mathbb C$ and spreads the finite continuation
data back to the stated characteristic-zero fields.

\subsection*{Foundational inputs}

We use the cone and contraction theorems, the klt existence and
relative big results of BCHM, surface minimal log discrepancy ACC,
lower-dimensional log MMP and dlt special termination, and Birkar's
rational complemented-pair existence theorem.  Their hypotheses are
specified at their applications.  On the topological side, Deligne's
mixed Hodge theory \cite{Deligne1971,Deligne1974}, the intersection
complexes of Goresky--MacPherson \cite{GM1983}, and the decomposition
theorem of Beilinson--Bernstein--Deligne--Gabber \cite{BBD1982}, with
Saito's Hodge-module refinement \cite{Saito1988,Saito1989}, supply
the weights and support bounds.  Formal functions, Grothendieck
existence and Artin approximation are used with the actual scheme
fibres.  Sections~\ref{sec:comparisons}--\ref{sec:lc} work over
$\mathbb C$; the field transfer applies to every program in the
theorem.  The termination proof is independent of the companion
results used for the endpoint consequences.
